% Contours de cartes à jouer à coins arrondis sur une page A4
%
% Format bridge (57 x 89 mm) : 11 cartes
%   2 colonnes de 3 cartes en portrait + 1 colonne de 5 cartes en paysage
%   largeur : 57 + 57 + 89 = 203 mm <= 210 mm
%   hauteur : 5 x 57 = 285 mm <= 297 mm
% Format poker (63 x 88 mm) : 9 cartes (grille 3 x 3)
%   Avec 63 + 88 = 151 mm on ne peut pas loger une colonne paysage
%   en plus de 2 colonnes portrait, donc 9 est le maximum pratique.
%
% Choisir le format ici :
\newif\ifbridge
%\bridgetrue      % 11 cartes bridge
\bridgefalse    % 9 cartes poker

\documentclass{article}
\usepackage[a4paper,margin=0mm]{geometry}
\usepackage{tikz}
\pagestyle{empty}

\begin{document}
\begin{tikzpicture}[remember picture,overlay,x=1mm,y=1mm,
    carte/.style={draw,line width=0.2pt,rounded corners=3mm}]
  \def\g{0}   % espace entre cartes (mm)
  \ifbridge
    \def\W{57}\def\H{89}
    % origine : coin inférieur gauche de la page
    \begin{scope}[shift={(current page.south west)}]
      \pgfmathsetmacro\xo{(210-(2*\W+\H+2*\g))/2}
      \pgfmathsetmacro\yp{(297-(3*\H+2*\g))/2}
      \pgfmathsetmacro\yl{(297-(5*\W+4*\g))/2}
      % 2 colonnes x 3 cartes en portrait
      \foreach \i in {0,1}
        \foreach \j in {0,1,2} {
          \pgfmathsetmacro\x{\xo+\i*(\W+\g)}
          \pgfmathsetmacro\y{\yp+\j*(\H+\g)}
          \draw[carte] (\x,\y) rectangle ++(\W,\H);
        }
      % 1 colonne x 5 cartes en paysage
      \foreach \j in {0,...,4} {
        \pgfmathsetmacro\x{\xo+2*(\W+\g)}
        \pgfmathsetmacro\y{\yl+\j*(\W+\g)}
        \draw[carte] (\x,\y) rectangle ++(\H,\W);
      }
    \end{scope}
  \else
    \def\W{63}\def\H{88}
    \begin{scope}[shift={(current page.south west)}]
      \pgfmathsetmacro\xo{(210-(3*\W+2*\g))/2}
      \pgfmathsetmacro\yo{(297-(3*\H+2*\g))/2}
      \foreach \i in {0,1,2}
        \foreach \j in {0,1,2} {
          \pgfmathsetmacro\x{\xo+\i*(\W+\g)}
          \pgfmathsetmacro\y{\yo+\j*(\H+\g)}
          \draw[carte] (\x,\y) rectangle ++(\W,\H);
        }
    \end{scope}
  \fi
\end{tikzpicture}
\null
\end{document}
